\documentclass[10pt,a4paper]{article}
\usepackage[utf8]{inputenc}
\usepackage[T1]{fontenc}
\usepackage[brazil]{babel}
\usepackage[margin=1.5cm]{geometry}
\usepackage{booktabs}
\usepackage{tabularx}
\usepackage{enumitem}
\usepackage{titlesec}
\usepackage{microtype}
\usepackage{hyperref}

\pagestyle{empty}

\titleformat{\section}{\large\bfseries}{}{0pt}{}[\titlerule]
\titlespacing*{\section}{0pt}{5pt}{3pt}

\setlist[itemize]{leftmargin=12pt,noitemsep,topsep=1pt,parsep=0pt,partopsep=0pt}
\setlist[enumerate]{leftmargin=12pt,noitemsep,topsep=1pt,parsep=0pt,partopsep=0pt}

\hypersetup{
    colorlinks=true,
    linkcolor=black,
    urlcolor=blue
}

\begin{document}

\begin{center}
    {\large \textbf{Relatório Técnico de Implementação -- Parte 2}}\\[2pt]
    \textbf{Disciplina:} DWCO8A -- Desenvolvimento Web \quad|\quad \textbf{Projeto:} Roteiros de Estudo\\
    \textbf{Autores:} Julia Romanetto dos Santos e Pedro Henrique Vogt Verolla \quad|\quad \textbf{Data:} 24 de outubro de 2026
\end{center}

\vspace{-4pt}

\section*{1. Instruções de Execução e Build Estático}
A aplicação foi construída com Next.js 16 (App Router) e React 19, utilizando CSS Modules nativos sem dependências externas de interface. Os comandos para execução e build são:
\begin{itemize}
    \item \textbf{Instalação de dependências:} execute \texttt{npm install} na raiz do projeto.
    \item \textbf{Ambiente de desenvolvimento:} execute \texttt{npm run dev} e acesse \url{http://localhost:3000}.
    \item \textbf{Exportação estática (Build):} execute \texttt{npm run build}. Com \texttt{output: "export"} configurado em \texttt{next.config.mjs}, o Next.js gera os arquivos HTML/CSS/JS compilados estaticamente na pasta \texttt{out/}.
    \item \textbf{Execução local do build estático:} execute \texttt{npx serve out} para subir um servidor HTTP local apontando para a pasta gerada, validando o funcionamento sem Node.js em tempo de execução.
\end{itemize}

\section*{2. API Externa Utilizada}
Foi integrada a API REST pública da Wikipédia em português através do endpoint:\\
\texttt{https://pt.wikipedia.org/api/rest\_v1/page/summary/\{tema\}}\\
A requisição não exige chave de autenticação (API pública do Wikimedia Foundation). No componente cliente \texttt{QuadroWikipedia}, a busca é disparada via \texttt{fetch} dentro de um \texttt{useEffect} usando o termo definido no roteiro. O ciclo de vida trata os três cenários de forma assíncrona e desacoplada:
\begin{itemize}
    \item \textbf{Carregamento:} exibe a mensagem textual \textit{"Carregando texto de apoio da Wikipédia..."} enquanto a requisição é processada.
    \item \textbf{Sucesso:} renderiza o resumo (\texttt{extract}) retornado e o link externo \textit{"Ler mais na Wikipédia"}.
    \item \textbf{Falha ou Erro de Rede:} caso \texttt{!res.ok} ou ocorra queda de conexão, exibe um aviso amigável explicando que o texto complementar não pôde ser carregado no momento, mas que o roteiro pode ser utilizado normalmente.
    \item \textbf{Não bloqueio:} a chamada ocorre no cliente e de modo isolado no rodapé da rota \texttt{/roteiros/[id]}, sem retardar a exibição do título, revisor, explicação principal ou questões.
\end{itemize}

\section*{3. Plano vs. Realizado}
A tabela abaixo compara o planejamento inicial definido na Parte 1 com a execução registrada nos commits do repositório:

\vspace{2pt}
\noindent
\begin{tabularx}{\textwidth}{c X c c c}
\toprule
\textbf{Etapa} & \textbf{Descrição do Escopo} & \textbf{Planejado} & \textbf{Realizado} & \textbf{Responsável} \\
\midrule
1 & Setup Next.js, ESLint, jsconfig e estrutura base & 19/10 & 20/10 & Julia \\
2 & Estruturação dos dados (\texttt{roteiros.json} e \texttt{professores.json}) & 20/10 & 20/10 & Julia \\
3 & Componente \texttt{Header} e navegação básica entre páginas & 20/10 & 20/10 & Julia \\
4 & Criação do componente \texttt{CardRoteiro} & 21/10 & 23/10 & Pedro \\
5 & Criação do componente \texttt{FiltroRoteiros} & 21/10 & 23/10 & Pedro \\
6 & Listagem dinâmica com estado e filtros na Home (\texttt{app/page.js}) & 22/10 & 23/10 & Pedro \\
7 & Rota dinâmica \texttt{/roteiros/[id]} e componente \texttt{DetalheRevisor} & 22/10 & 23/10 & Pedro \\
8 & Componente \texttt{ItemQuestao} com controle de gabarito (\texttt{useState}) & 23/10 & 23/10 & Pedro \\
9 & Componente \texttt{QuadroWikipedia} consumindo API REST externa & 23/10 & 23/10 & Pedro \\
10 & Rota \texttt{/solicitar} e \texttt{FormularioSolicitacao} com inputs controlados & 24/10 & 23/10 & Pedro \\
11 & Validação de regras de negócio no formulário e feedback visual & 24/10 & 23/10 & Pedro \\
12 & Ajuste do build estático (\texttt{generateStaticParams} com id \texttt{'999'}) & 25/10 & 23/10 & Pedro \\
13 & Testes gerais, verificação do build e elaboração deste relatório & 25/10 & 24/10 & Pedro e Julia \\
\bottomrule
\end{tabularx}

\vspace{1pt}

\section*{4. Mudanças em relação ao Plano e ao Protótipo}
\begin{itemize}
    \item \textbf{Adiantamento do Cronograma:} as etapas de 4 a 12, planejadas para ocorrerem entre 21/10 e 25/10, foram concentradas e concluídas em 23/10 por Pedro. A antecipação permitiu validar precocemente o build estático e os testes da API, liberando o dia 24/10 exclusivamente para revisão e relatório.
    \item \textbf{Opção "Todos" no Filtro de Anos:} o protótipo inicial previa apenas opções rígidas de 6º a 9º ano. Adicionou-se a opção \texttt{"Todos"} no \texttt{<select>} de anos para permitir a visualização irrestrita dos roteiros por disciplina sem forçar a escolha de uma série escolar específica.
    \item \textbf{Tratamento de Roteiro Inexistente sem \texttt{notFound()}:} para atender à exportação estática com a rota de teste \texttt{\{ id: '999' \}}, substituiu-se o redirecionamento 404 por uma renderização amigável na própria página com mensagem \textit{"Roteiro não encontrado"} e link de retorno para \texttt{/}.
    \item \textbf{Validação Personalizada no Formulário (\texttt{noValidate}):} utilizou-se \texttt{noValidate} no elemento \texttt{<form>} para desabilitar balões nativos do navegador e permitir o acúmulo de mensagens claras em uma caixa de feedback em vermelho, conforme os requisitos pedagógicos.
    \item \textbf{Ajuste de Hidratação (\texttt{suppressHydrationWarning}):} adicionado ao \texttt{<body>} em \texttt{app/layout.js} para neutralizar discrepâncias de atributos injetados por extensões do navegador (como o ColorZilla) em ambiente de desenvolvimento.
\end{itemize}

\end{document}
